\section{Results}
\label{sec:results}

\subsection{Semantic Similarity Analysis}
\label{sec:res_semantic}

\Tabref{tab:semantic_overall} reports top-level metrics comparing \semsim predictions
against \gptmini selections.  Cosine similarity achieves 40\% top-1 accuracy---60\%
better than the 25\% random baseline.  The mean rank of the LLM's chosen tool is 2.09
out of 4 (chance rank 2.5), confirming a systematic bias toward higher-similarity tools.
However, the mean cosine similarity of the \emph{selected} tool (0.219) is meaningfully
lower than that of the \emph{top-similar} tool (0.257), indicating that the LLM does not
simply select the highest-cosine-similarity option.

Mixed-category examples (where tools differ in domain) reach 60\% top-1 accuracy,
compared to 37.5\% for same-category examples.  This gap is expected: when tools span
distinct domains, semantic differences are more pronounced and cosine similarity is more
discriminative.

\para{Per-category analysis.}  \Tabref{tab:per_category} breaks down semantic accuracy
and the position-0 selection rate by tool category.  Calculator queries achieve the
highest semantic accuracy (66.7\%) because mathematical vocabulary (``solve'', ``compute'',
``equation'') aligns well with tool descriptions containing ``computation'' and
``mathematical expressions''.  Translation queries show the lowest accuracy (22.2\%)
because all translation tools share similar descriptions mentioning ``languages'', ``text'',
and ``translate'', making cosine similarity nearly uninformative for within-category
discrimination.  The high position-0 rate in the main experiment (57.8\% overall) is
consistent with a correlation between description quality and list position in our
dataset; the controlled rotation experiment below separates this confound.

\begin{table}[t]
\centering
\caption{Semantic similarity results for the main selection experiment (45 LLM selections).
         Top-1 accuracy measures how often the highest-similarity tool matches the LLM's
         selection.  Chance for all accuracy metrics is 0.25.}
\label{tab:semantic_overall}
\begin{tabular}{@{}lcc@{}}
\toprule
\textbf{Metric} & \textbf{Value} & \textbf{Chance} \\
\midrule
Top-1 accuracy (overall)          & \textbf{0.400} & 0.250 \\
Top-1 accuracy (same-category)    & 0.375          & 0.250 \\
Top-1 accuracy (mixed-category)   & \textbf{0.600} & 0.250 \\
Mean rank of LLM's selection      & 2.09           & 2.50  \\
Mean cosim (selected tool)        & 0.219          & ---   \\
Mean cosim (top-similar tool)     & 0.257          & ---   \\
Position-0 selection rate         & 0.578          & 0.250 \\
\bottomrule
\end{tabular}
\end{table}

\begin{table}[t]
\centering
\caption{Per-category semantic top-1 accuracy and position-0 selection rate.
         Calculator queries benefit from distinctive mathematical vocabulary;
         translation queries suffer from high within-category description similarity.}
\label{tab:per_category}
\begin{tabular}{@{}lccc@{}}
\toprule
\textbf{Category} & \textbf{$n$} & \textbf{Semantic Top-1} & \textbf{Position-0 Rate} \\
\midrule
Calculator      & 9 & \textbf{0.667} & 0.667 \\
Weather         & 9 & 0.556          & 0.556 \\
Code execution  & 8 & 0.375          & 0.500 \\
Search          & 8 & 0.250          & 0.500 \\
Translation     & 9 & 0.222          & 0.667 \\
\bottomrule
\end{tabular}
\end{table}

\subsection{Positional Bias Analysis}
\label{sec:res_pos}

\Tabref{tab:positional_bias} shows the empirical selection rate at each position in the
controlled rotation experiment.  Position-0 is selected 32.5\% of the time---above chance
(25\%) but not significantly so ($\chi^2 = 1.80$, $p = 0.615$, $n = 40$).  The total
variation distance from the uniform distribution is only 0.075.  These results sharply
contrast with the 57.8\% position-0 rate observed in the uncontrolled main experiment,
where tools were not systematically rotated.  The most striking finding is selection
stability: the same tool is selected across all four orderings in 90\% of base examples,
demonstrating that the model tracks semantic relevance regardless of list position.

\begin{table}[t]
\centering
\caption{Position preference in the rotation experiment (10 base examples $\times$ 4
         rotations $= 40$ variants).  No position is selected significantly more than
         chance (25\%).  Selection stability: 90\% of base examples produce the same
         winner across all four rotations.}
\label{tab:positional_bias}
\begin{tabular}{@{}lcc@{}}
\toprule
\textbf{Position} & \textbf{Observed Rate} & \textbf{Expected Rate} \\
\midrule
Position 0 & 0.325 & 0.250 \\
Position 1 & 0.250 & 0.250 \\
Position 2 & 0.250 & 0.250 \\
Position 3 & 0.175 & 0.250 \\
\midrule
\multicolumn{3}{@{}l@{}}{\small $\chi^2 = 1.80$, $p = 0.615$; Selection stability: \textbf{0.900}} \\
\bottomrule
\end{tabular}
\end{table}

\subsection{Description Perturbation}
\label{sec:res_perturb}

\Tabref{tab:perturbation} reports the selection under original, generic, and adversarial
description conditions.  The key finding is that \emph{generic} descriptions (neutral,
semantically uninformative) do not change selection in any of the three cases, suggesting
the model retains a prior from the tool name itself.  However, \emph{adversarial}
descriptions (descriptions swapped between tools) change selection in 2 out of 3 cases
(67\%).  In both changed cases, the model follows the description rather than the tool
name, directly confirming that description content causally drives selection.

The one failure case---the translation query---is informative: even with an adversarial
swap, the model continues selecting \texttt{google\_translate\_api}, suggesting that the
name ``Google Translate'' carries a strong semantic prior that overrides the misleading
description.  This hints at a secondary signal path from tool \emph{name} to selection
that does not go through description content.

\begin{table}[t]
\centering
\caption{Description perturbation results across three query types and three conditions.
         ``Changed?'' indicates whether the selected tool differs from the original
         condition.  Adversarial descriptions swap tool descriptions; generic
         descriptions replace all descriptions with a neutral placeholder.}
\label{tab:perturbation}
\begin{tabular}{@{}llll@{}}
\toprule
\textbf{Query Type} & \textbf{Condition} & \textbf{Selected Tool} & \textbf{Changed?} \\
\midrule
Weather (London)    & Original    & openweathermap\_api   & ---        \\
Weather (London)    & Generic     & openweathermap\_api   & No         \\
Weather (London)    & Adversarial & wolfram\_alpha\_api   & \textbf{Yes} \\
\midrule
Calculation ($25\times48$) & Original    & wolfram\_alpha\_api   & ---        \\
Calculation ($25\times48$) & Adversarial & openweathermap\_api  & \textbf{Yes} \\
\midrule
Translation (Spanish) & Original    & google\_translate\_api & ---       \\
Translation (Spanish) & Adversarial & google\_translate\_api & No        \\
\midrule
\multicolumn{4}{@{}l@{}}{\small Adversarial change rate: 2/3 (67\%); Generic change rate: 0/1 (0\%)} \\
\bottomrule
\end{tabular}
\end{table}

\subsection{Layer-wise Probing Results}
\label{sec:res_probing}

\para{Task 1: Tool category probe.}
\Tabref{tab:probe_category} shows that tool category is perfectly decodable from
\emph{every} layer of \gpttwo using mean-pooled representations.  Probe accuracy is
100\% in layers 1--8 and 99.8\% in layers 9--12, all far above the 20\% chance baseline.
This result suggests that semantic category information is primarily encoded at the
token--embedding level and is immediately available to all subsequent layers.

\para{Task 2: Correct tool position probe (mean pooling).}
\Tabref{tab:probe_position_mean} shows probe accuracy for predicting which tool slot
is semantically correct, using mean-pooled representations.  In sharp contrast to the
category results, accuracy rises progressively: from 35.0\% at layer~1 (10 points above
chance) to a peak of 68.0\% at layer~9, before declining to 45.0\% at layer~12.  This
gradient---visualized in \figref{fig:probing}---is the paper's central mechanistic finding.
Layers 7--9 are the critical window for semantic matching computation.

\para{Task 3: Correct tool position probe (last-token).}
When probing the representation at the \emph{last token} position (the generation
point in autoregressive decoding), accuracy is uniformly high across all layers: 72.3\%
in early layers, 73.3\% in middle layers, and 73.0\% in late layers, with a peak of
81.0\% at layer~8 or~10.  This high early accuracy---compared to only 35\% for mean
pooling at layer~1---suggests that the last-token representation immediately aggregates
the query context relevant to selection, consistent with its role as the autoregressive
generation site.

\begin{table}[t]
\centering
\caption{Tool category probe accuracy by layer group (\gpttwo, mean pooling, 5-fold CV).
         Category is perfectly decodable from all layers; chance is 0.20.}
\label{tab:probe_category}
\begin{tabular}{@{}lcc@{}}
\toprule
\textbf{Layer Group} & \textbf{Mean Accuracy} & \textbf{Chance} \\
\midrule
Early (1--4)  & \textbf{1.000} & 0.20 \\
Middle (5--8) & \textbf{1.000} & 0.20 \\
Late (9--12)  & 0.998          & 0.20 \\
\bottomrule
\end{tabular}
\end{table}

\begin{table}[t]
\centering
\caption{Correct tool position probe accuracy by layer (\gpttwo, mean pooling, 5-fold CV).
         Accuracy rises from 35\% at layer~1 to 68\% at layer~9, then declines.
         Chance is 0.25.}
\label{tab:probe_position_mean}
\renewcommand{\arraystretch}{1.15}
\begin{tabular}{@{}lccc@{}}
\toprule
\textbf{Layer} & \textbf{Accuracy} & \textbf{Std} & \textbf{vs.\ Chance} \\
\midrule
1  & 0.350 & 0.077 & $+$10.0\% \\
2  & 0.360 & 0.058 & $+$11.0\% \\
3  & 0.410 & 0.049 & $+$16.0\% \\
4  & 0.440 & 0.058 & $+$19.0\% \\
5  & 0.490 & 0.097 & $+$24.0\% \\
6  & 0.560 & 0.066 & $+$31.0\% \\
7  & 0.600 & 0.105 & $+$35.0\% \\
8  & 0.600 & 0.089 & $+$35.0\% \\
\textbf{9}  & \textbf{0.680} & \textbf{0.081} & $\mathbf{+}$\textbf{43.0\%} \\
10 & 0.680 & 0.075 & $+$43.0\% \\
11 & 0.630 & 0.081 & $+$38.0\% \\
12 & 0.450 & 0.130 & $+$20.0\% \\
\bottomrule
\end{tabular}
\end{table}

\begin{figure}[t]
    \centering
    \includegraphics[width=0.95\linewidth]{fig5_layer_probing.png}
    \caption{Layer-wise probe accuracy for correct tool position in \gpttwo (mean pooling).
             Accuracy rises progressively from 35\% at layer~1 to a peak of 68\% at
             layer~9, then declines toward the final layer.  The dashed line marks the
             25\% chance baseline.  This gradient indicates that semantic matching is
             a multi-layer computation concentrated in the middle-to-late layers.}
    \label{fig:probing}
\end{figure}

\begin{figure}[t]
    \centering
    \begin{subfigure}[b]{0.48\textwidth}
        \includegraphics[width=\textwidth]{fig1_semantic_accuracy.png}
        \caption{Semantic similarity top-1 accuracy vs.\ baselines across categories.}
        \label{fig:semantic}
    \end{subfigure}
    \hfill
    \begin{subfigure}[b]{0.48\textwidth}
        \includegraphics[width=\textwidth]{fig4_perturbation.png}
        \caption{Selection changes under original, generic, and adversarial descriptions.}
        \label{fig:perturbation}
    \end{subfigure}
    \caption{Behavioral results.  \captiona{}~Semantic similarity is above chance for all
             categories but falls short for translation (within-category descriptions are
             nearly identical).  \captionb{}~Adversarial perturbation causes the model to
             follow description content over tool name in 67\% of cases; generic
             descriptions leave selection unchanged.}
    \label{fig:behavioral}
\end{figure}
